\subsection{N-021: an all-degree finite reduction for cancelling-pair minimization}
\label{sec:dimension-degree-cutoff}

\paragraph{Scope.}
This is a finite-search reduction for the remaining dimension-ring
question.  It does not compute a global minimum or identify a natural
Hecke-category selection.  Unlike the earlier polynomial-parameter
searches, it applies to all Laurent dimension tables on the prescribed
basis, with no assumed cell formula or degree window.  The budget
counts all cancelling pairs, not just those in the bottom cell.

Let $c_{ij}^k\in\mathbb Z[q,q^{-1}]$ be an associative unital table
on a finite ordered basis with designated unit.  Write $|c|$ for
coefficientwise absolute value, and let $d_0$ bound the absolute
values of all exponents appearing in $c$.  A dimension table $d$
means an associative unital table on the same free module with
\[
 d_{ij}^k\in\mathbb N[q,q^{-1}],\qquad
 d_{ij}^k\ge |c_{ij}^k|,\qquad d_{ij}^k\equiv c_{ij}^k\pmod2.
 \tag{N-115}
\]
Its cancelling-pair score is
$B(d)=\frac12\sum_{i,j,k}(d_{ij}^k(1)-|c_{ij}^k|(1))$.
Prescribed zero slots may be retained, including the unit constraints.
One may additionally require bar symmetry $d(q)=d(q^{-1})$ and
specified label-permutation or opposite-table symmetries.

\paragraph{The finite bound.}
If a dimension table of score $b\le B$ exists, where $B\ge1$,
then one exists with the same score $b$ and all exponents bounded by
\[
                 K(B,d_0)=B(d_0+1)4^B.
 \tag{N-116}
\]
The zero-score case is just the single table $|c|$ and is checked
directly.  The bound preserves the optional symmetries just listed
and every prescribed zero slot.  It need not preserve unimodality;
that case has the different, much smaller bound below.

The conclusion bounds a replacement table, not every table of
the same score.  For example, for the true two-label ring
$x^2=1+x$, each table
$x^2=1+(1+2q^N+2q^{-N})x$ has score2, for arbitrarily large $N$.
It can be replaced by $x^2=1+5x$, with the same score, unit and
bar symmetry.

\paragraph{Encode extra pairs by exponent variables.}
Every table satisfying (N-115) has an expression
\[
 d_{ij}^k=|c_{ij}^k|+
                2\sum_{t\in T_{ij}^k}q^{x_t},
 \qquad \sum_{i,j,k}|T_{ij}^k|=b,
 \quad x_t\in\mathbb Z.
 \tag{N-117}
\]
Occurrences are distinguished even when their exponents coincide.
Fix the allocation of these $b$ pair occurrences to slots from
an existing solution.  The unit rows and columns have no extra
pairs, because their coefficients are already the required zeros
and ones.  If a slot is prescribed to stay zero, no pair is
allocated to it.  Changing only the exponents will preserve these
conditions and the exact score, including when occurrences merge.

Expand every nonnegative integer coefficient into that many
monomial occurrences; each pair occurrence in (N-117) contributes
two copies with the same exponent variable.  For any associativity
identity
\[
 \sum_\ell d_{ij}^{\ell}d_{\ell k}^{m}
       =\sum_\ell d_{jk}^{\ell}d_{i\ell}^{m},
 \tag{N-118}
\]
the existing solution gives a bijection between the monomial
occurrences on its two sides, matching equal exponents.  Fix one
such bijection for every identity.  Each matched equality is an
integer affine equation in the $x_t$: it compares sums of at most
two exponents on each side.  Thus the resulting integer system
$Ax=b'$ has
\[
       \sum_t|A_{st}|\le4,
                  \qquad |b'_s|\le4d_0.
 \tag{N-119}
\]
The matching depends on the original solution; it is fixed before
solving this affine system.  The argument does not claim that one
arbitrary matching works for every table.  Conversely, any integer
solution of this fixed system makes every matched pair of monomials
equal, and hence makes the entire table associative.

Bar symmetry is retained by also fixing monomial bijections
between each table entry and its bar.  These give equations with
row norm at most2 and constants of size at most $2d_0$.
When $|c|$ itself is bar fixed, this can equivalently be done on
the extra pair occurrences: pair $x_t$ with $-x_{t'}$, and use
$x_t=0$ for occurrences fixed by that pairing.  Label permutations
and opposite-table symmetries are imposed by the same matching
construction, with row norm at most2.  They do not increase (N-119).
The argument is in $q=v^2$, so every resulting Laurent exponent
in $v$ remains even.  If a different common exponent lattice
$q^h$ is prescribed and contains the input, run the argument in
that lattice variable; no forbidden exponent is introduced.

\paragraph{Small integer solution lemma, with its full proof.}
Consider any solvable integer affine system in $b$ variables
satisfying (N-119), and let its rank be $r$.  If $r=0$, consistency
makes all right sides zero and the zero vector works.  Otherwise
choose $r$ independent rows $I$ and $r$ columns $P$ whose square
submatrix $H$ is nonsingular.  Put $\Delta=|\det H|$ and let $N$
be the other columns.  Hadamard's bound and the row norm give
\[
 \Delta\le4^r,
 \qquad |\operatorname{adj}(H)_{ij}|\le4^{r-1}.
 \tag{N-120}
\]
Start with any integer solution $x$.  Replace each nonpivot
coordinate by its residue $x'_t\in\{0,\ldots,\Delta-1\}$ modulo
$\Delta$.  Define the pivot coordinates by
\[
       x'_P=H^{-1}(b'_I-A_{I,N}x'_N).
 \tag{N-121}
\]
They are integers: the change in each nonpivot coordinate is a
multiple of $\Delta$, and
$\Delta H^{-1}=\pm\operatorname{adj}(H)$ is an integer matrix.
The independent equations are satisfied by construction.  Every
other row of $A$ is a rational linear combination of rows $I$;
because the original system has a solution, its right side is
the same combination of $b'_I$.  Consequently all dependent
equations also remain satisfied.

For the numerical bound, each coordinate of
$b'_I-A_{I,N}x'_N$ has absolute value at most
$4d_0+4(\Delta-1)$.  Applying the adjugate estimate gives
\[
 \begin{split}
 |x'_t|&\le r4^{r-1}\bigl(4d_0/\Delta+4\bigr)\\
       &\le r(d_0+1)4^r
          \le B(d_0+1)4^B\qquad(t\in P).
 \end{split}
 \tag{N-122}
\]
For $t\in N$, $0\le x'_t<\Delta\le4^r$ gives the same final
bound.  Reinsert these exponents in (N-117).  The fixed matching
identities prove associativity and the chosen symmetries; (N-115),
the exact score and the unit are preserved by construction.
Since $K(B,d_0)\ge d_0$, all original exponents also fit in
the box.  This proves (N-116).

\paragraph{A sharper bound for symmetric unimodal dimensions.}
Let $M=\max_{i,j,k}|c_{ij}^k|(1)$.  If every nonzero dimension
entry is symmetric and unimodal, its support is an interval
$[-a,a]$ with a positive integer coefficient in every position.
Its coefficient sum is therefore at least $2a+1$.  A table of
score at most $B$ satisfies, entry by entry,
$d_{ij}^k(1)\le M+2B$, so every exponent obeys
\[
                    |e|\le
             \left\lfloor\frac{M+2B-1}{2}\right\rfloor.
 \tag{N-123}
\]
This includes entries whose true trace constant is zero; if their
dimension entry is nonzero, the same symmetry and interval argument
applies.  Zero entries need no bound.  Unlike (N-116), (N-123)
holds for every table in the stated class, without compression.

\paragraph{Application and honest computational boundary.}
For the132 nonzero rank-three constants supplied with
$E^{\rm opp}(3)$, direct exact extraction gives
\[
 d_0=5,\qquad
 \sum_{i,j,k}|c_{ij}^k|(1)=4443,\qquad M=576.
 \tag{N-124}
\]
The sum includes19 unit constants, agreeing with N-005's4424
nonunit total.  Its previously verified examples have scores
7276 in general and9040 with all dimensions unimodal.
Consequently some global minimizer among all Laurent dimension
tables lies in the completely explicit degree box
\[
             |e|\le43656\,4^{7276}.
 \tag{N-125}
\]
The degree bound has14568 binary digits.  Optional bar and
table-label symmetries can be retained if one restricts the
minimization to those classes and uses an existing upper-bound
example in that class.  Every symmetric unimodal minimizer
with score at most9040 already lies in the much smaller box
\[
                         |e|\le9327.
 \tag{N-126}
\]
For any such box and budget there are finitely many allocations,
exponents and nonnegative integer coefficients.  Direct sparse
polynomial multiplication checks every resulting table.  Enumerating
scores from zero through the stated upper bound is therefore an
exact terminating decision procedure for the global minimum.
No such exhaustive optimization is claimed to have been performed.
In particular neither7276 nor9040 is promoted to a global minimum,
and this result gives no geometric or canonical selection principle.
The theorem applies directly to the integral Laurent trace table
at rank three.  N-019's rational-coefficient alternative is not
automatically an input satisfying (N-115), and no conclusion for
an undefined $E(4)$ is obtained by applying this bound.

One further complexity consequence is precise.  Given a sparse
integer Laurent trace table and a score budget written in unary,
existence of a dimension table is in NP, including the optional
symmetries above.  A certificate is the at most $B$ extra pair
occurrences, their slots, and their bounded binary exponents.
Equation (N-116) makes its length polynomial in $B$ and the input
length.  Associativity is verified by sparse polynomial products;
coefficient magnitudes are manipulated in binary.  Expanding
coefficients into individual monomials was used only in the
existence proof of the bound, not in this verifier.  This is not
a polynomial-time optimization algorithm, and does not claim
NP membership with a binary-encoded pair budget.

\paragraph{Evidence and standard-method credit.}
The determinant and small-integer-solution method is standard;
compare Papadimitriou (1981), Lemma1, printed p.~766, and its
integer-programming small-witness argument.  The particular
monomial-matching reduction and bounds here are proved above,
without importing a general integer-programming bound.
\nolinkurl{dimension_exponent_compression.py} implements the
integer-system compression, checks dependent rows, and supplies
an explicit two-label example with arbitrarily remote exponents.
It also extracts (N-124) exactly from the supplied constant file.
The script is an implementation check of the finite mechanism;
the proof, rather than a sample, establishes all-degree validity.
